\chapter{EZDrone: La estación de tierra integrada}
Los capítulos anteriores han recorrido el camino desde la arquitectura original basada en HTTP hasta la arquitectura híbrida MQTT/WebRTC descrita en el capítulo 7. Este capítulo presenta finalmente EZDrone, el sistema resultante de integrar esa arquitectura de comunicaciones con una aplicación de control completa. Mientras que el Capítulo 7 se centra en cómo viajan los datos, este se centra más en qué ofrece la aplicación al operador, la gestión centralizada del estado, los tres modos de control, el soporte de tres tipos distintos de drones, la plataforma de misiones, la visualización en tiempo real y el registro de los vuelos. El capítulo cierra con una demostración integral del sistema en funcionamiento.

\section{Visión general de EZDrone}
EZDrone es una GCS que se ejecuta en el propio navegador. A diferencia de las soluciones de escritorio clásicas como Mision Planner o QGroundControl, no requiere instalación previa, la aplicación Flutter Web se sirve como un conjunto de ficheros estáticos a los que el operador accede desde cualquier dispositivo con cualquier navegador. Esta decisión de diseño es la que habilita los modos de control por voz e inercial, que se apoyan en sensores y APIs propias del navegador del dispositivo.

El sistema se asienta sobre la arquitectura híbrida del Capítulo 7, que reparte cada flujo de información por el protocolo más adecuado a sus características. Los comandos discretos y las confirmaciones de estado viajan por MQTT sobre WebSocket seguro, los flujos de alta frecuencia viajan por canales WebRTC peer-to-peer, que aportan la baja latencia y el ancho de banda que MQTT no puede ofrecer. 

Sobre esta base, EZDrone añade una capa funcional que es el objeto de este capítulo, un único punto de gestión del estado (el DronProvider), tres modos de control alternativos (clásico, por voz e inercial), soporte para drones (ArduPilot reales, el simulador SITL, y el DJI Tello), planificación y ejecución de misiones autónomas, visualización cartográfica y de vídeo en tiempo real, y un sistema de registro y galería multimedia persistente en el propio navegador.

\section{Gestión de estado: el DronProvider y la máquina de estados}
Toda la lógica de estado de la aplicación se concentra en una única clase, el DronProvider. Los widgets nunca almacenan el estado de vuelo, sino que se limitan a leerlo para redibujarse (\texttt{context.watch}) o a invocar acciones sobre él (\texttt{context.read}). Esta centralización es aposta, ya que garantiza que las distintas pantallas observen siempre una única fuente, y queden desacopladas de la lógica de transporte (MQTT y WebRTC).

El Provider mantiene tres tipos de estado. En primer lugar, un conjunto de enumeraciones que fijan la configuración de la sesión, \texttt{ControlMode} (clásico, voz, IMU), \texttt{DroneConnectionMode} (arduPilot, sitl, tello) y \texttt{DetectionMode} (all, person, none). En segundo lugar, una colección de variables booleanas que dirigen qué botones se habilitan y que muestra la interfaz (\texttt{isConnected, isArmed, isConfigValid, isMissionMode,} entre otras). Y en tercer lugar, las variables de telemetría (\texttt{currentLat, currentLon, currentAlt, currentBat, currentHeading...}), que se actualizan a partir del canal de telemetría WebRTC.

Estas variables no son independientes entre sí, sino que están ligadas por una máquina de estados que reproduce el ciclo de un dron real, (\texttt{deconectado - conectado - armado - volando - aterrizado}). Cada transición dispara una confirmación enviada por la estación de tierra, no por el propio comando. Es decir, al pulsar ARM, la aplicación no se considera armada hasta que recibe el mensaje \texttt{armed} de vuelta. Hasta entonces, permanece en un estado de espera. Esto evita que la interfazz se adelante al estado físico del dron.  

(tabla de botones y booleanos??????)

\section{Soporte multi-dron}
Una de las aportaciones de EZDrone respecto a la arquitectura original es el soporte de tres destinos de vuelo distintos, seleccionables desde la pantalla de configuración inicial sin cambios en el cliente. La elección se publica al conectar mediante el topic \texttt{setMode} y determina como la estación tierra establece el enlace.

Los modos \texttt{ardupilot} y \texttt{sitl} comparten toda la lógica de pilotaje, los mismos modos y funciones, basados en la librería dronLink, que implementa el protocolo MAVLink, y descompone las diferentes acciones (armar, despegar, ir a un punto, telemetría etc.) en diferentes módulos propios. La única diferencia entre ambos es el extremo de conexión, un autopilot Pixhawk en el caso de dron real, o el proceso del simulador SITL en el caso del simulado. Esta equivalencia en general, es la que permite desarrollar y validar el sistema sin necesidad de probarlo en el hardware real.

El modo \texttt{tello} recurre a una librería distinta, la de TelloLink, ya que DJI no usa MAVLink, sino su propio SDK sobre WI-FI. La pantalla del Tello en si también es una implementación separada (\texttt{TelloFlightScreen}), diseñada intentando replicar la propia pantalla de control de DJI en torno a los modos especiales que ofrece el tello y que se describen en el apartado 8.6.

\section{Modos de control (ArduPilot / SITL)}
EZDrone ofrece tres formas alternativas de pilotar el dron en los modos ArduPilot y SITL. Todas producen el mismo tipo de salida (valores normalizados de cuatro ejes enviados por el DataChannel del joystick), pero difieren radicalmente en la interfaz de entrada. El modo se elige en la pantalla de configuración y determina a qué pantalla de vuelo se accede.

\subsection{Modo clásico (FlightScreen)}
La pantalla clásica (FlightScreen) reproduce el control de un mando de rc mediante dos joysticks virtuales y un conjunto de instrumentos. En modo apaisado, la pantalla se divide en tres columnas, el panel izquierdo muestra instrumentos numéricos (altitud, velocidad, latitud y longitud), el panel central contiene el mapa satélite o el vídeo FPV con los dos joysticks superpuestos y el panel derecho añade la rosa de los vientos animada según la orientación, el propio heading y el estado del vuelo. En modo vertical, los instrumentos se compactan en una sola fila horizontal en la parte superior.

El stick izquierdo lleva el thrust y el yaw, mientras que el derecho el pitch y roll. Los ejes del stick izquierdo se atenúan por un factor de 0.35 para suavizar las entradas (ya que es un movimiento muy brusco). La barra de acción inferior contiene los botones ARM, TAKEOFF, LAND, RTL y DISCONNECT, cuya habilitación depende directamente de los booleanos del Provider descritos en el apartado 8.2.

Sobre el vídeo en directo, se superponen varios controles. Un dropdown para seleccionar el modo de detección YOLO (Todo/personas/nada), botones de captura de foto y grabación de vídeo, un selector de cámara (dron/portátil con ET) y una barra de zoom (de 1x a 5x) que se orienta verticalmente en modo portrait y horizontalmente en modo landscape. Cuando el vídeo se está grabando, aparece el badge de \texttt{REC} en la esquina durante toda la grabación